\section{The product}\label{sec:product}

\paragraph{Who uses it.} Road haulage in Italy is dominated by very small firms: the typical
\emph{autotrasportatore} runs a handful of trucks. Diesel is one of their largest costs, typically around
30\% of operating costs, and their margins are a few percent of revenue. A ten-truck firm burns about
25\,000 litres a month; in our data the wholesale price rose by 32\% in a single month in March 2022
(about 5\,000~EUR on a month's fuel bill) and by 60\% in March 2026. These firms do not trade futures, but
they can ask their fuel supplier or their bank for a fixed price for next month. The question they cannot
answer is \emph{when} fixing is worth its cost (we assume a 1\% margin over the futures price).

\paragraph{The client.} We pitch FuelGuard to a \textbf{bank with SME loans to transport firms}. Fuel
shocks are a credit-risk driver for these borrowers, and the bank can sell fixed-price fuel or simple
swaps. Each month FuelGuard gives the bank, for each client, two numbers it can pass on through its
relationship managers or app:
\begin{enumerate}
\item a \textbf{hedge ratio} (0\%, 50\% or 100\% of next month's diesel), and
\item a \textbf{95\% worst case}: ``next month, there is a 5\% chance that your diesel will cost more than
X\% above the price you can fix today''.
\end{enumerate}

\paragraph{The prediction.} Both numbers come from one machine-learning task: a \textbf{regression of next
month's realised volatility} of diesel futures. Our first design tried to predict the price itself, or
whether it would jump by more than 5\%. Out of sample, the size of next month's price move could not be
predicted beyond the futures price, and the modest skill in predicting jumps came almost entirely from
implied volatility (Appendix~\ref{app:v1}). Following the literature, which finds that volatility is
persistent and predictable while price changes are not \citep{corsi2009,christensen2023,degiannakis2017},
FuelGuard~2.0 forecasts \emph{how much} the price will move, not in which direction. For a hedger this is
the relevant quantity: the risk of a costly month grows with volatility.

\paragraph{What it replaces (benchmarks).} The haulier's current practice (never fix, always fix, fix half)
for the decision; and, for the forecast, the three models a risk desk would use today: last month's
volatility (random walk), the HAR model of \citet{corsi2009}, and HAR augmented with the option-implied
volatility (HAR-IV), the market's own forecast. HAR-IV is the bar the machine-learning models must clear.
